\chapter{Заключение}
\label{ch:conclusion}

Разработка системы рекомендаций для фильмов стала значительным шагом в применении современных технологий машинного обучения и веб-разработки для решения одной из актуальных задач цифровой эпохи — персонализации контента. В ходе выполнения проекта была создана платформа, объединяющая функциональность, безопасность и удобство использования, что соответствует требованиям современных пользователей и бизнесов.

Основные результаты проекта включают:
\begin{itemize}
    \item Разработку надёжного и масштабируемого бэкенда, который обеспечивает хранение и обработку данных о фильмах,
    \item Интеграцию продвинутых алгоритмов машинного обучения для предоставления персонализированных рекомендаций, учитывающих предпочтения пользователей и контекст текущего просмотра,
    \item Создание удобного и интуитивно понятного интерфейса для пользователей, который предоставляет функционал поиска, просмотра информации о фильмах, а также оценки и взаимодействия с контентом,
    \item Реализацию защитных механизмов для сохранения конфиденциальности данных пользователей и обеспечения их безопасности.
\end{itemize}

Проект был выполнен командой специалистов с использованием передовых методологий, таких как Agile, что позволило эффективно распределить задачи и ресурсы, а также обеспечить адаптивность к изменениям. Применение инструментов, таких как Jira, Confluence, Docker, способствовало ускорению работы, поддержанию высокого качества разработки и обеспечению стабильного функционирования системы.

Кроме того, особое внимание было уделено тестированию системы, что позволило обнаружить и исправить возможные ошибки на ранних этапах разработки, повысив надёжность и производительность платформы.  

Проект не только достиг поставленных целей, но и заложил основу для дальнейшего развития. Возможными направлениями улучшений могут быть:
\begin{itemize}
    \item Расширение функционала системы рекомендаций, включая внедрение гибридных моделей машинного обучения,
    \item Интеграция с внешними платформами для получения дополнительных данных о фильмах,
    \item Оптимизация архитектуры для ещё большей масштабируемости и производительности. Интеграция Airflow для автоматического обучения модели на новых данных.
\end{itemize}

В целом, проект демонстрирует, как современные технологии и эффективное управление разработкой могут быть использованы для решения задач любой сложности. Итоги работы команды являются вкладом в развитие технологий персонализации контента, и платформа обладает потенциалом для дальнейшего коммерческого использования и масштабирования.

\endinput